%%% Frontiers in Artificial Intelligence — Original Research
%%% Uses the Frontiers Harvard (author-date) template (FrontiersinHarvard.cls).
%%% Compile: pdflatex; bibtex; pdflatex; pdflatex  (Frontiers-Harvard.bst + references.bib)

\documentclass[utf8]{FrontiersinHarvard}
\usepackage{url,hyperref,lineno,microtype,subcaption}
\usepackage{booktabs}
\usepackage{array}
\usepackage{calc}
\usepackage{longtable}
\usepackage[onehalfspacing]{setspace}
\linenumbers

\def\keyFont{\fontsize{8}{11}\helveticabold }
\def\firstAuthorLast{Chincholikar {and} Chawla}
\def\Authors{Samir Chincholikar\,$^{1}$, Robin Chawla\,$^{2,*}$}
\def\Address{$^{1}$Independent Researcher, New York, NY, United States \\ $^{2}$Independent Researcher, India}
\def\corrAuthor{Robin Chawla}
\def\corrEmail{robin.chawla.cse14@iitbhu.ac.in}

\begin{document}
\onecolumn
\firstpage{1}

\title[The AI Incident Disclosure Shadow]{The AI Incident Disclosure Shadow: An Incident-Level Study Matching AIID-Recorded AI Failures to U.S.-Listed Companies' Investor Disclosures, 2019--2026}

\author[\firstAuthorLast ]{\Authors}
\address{}
\correspondance{}
\extraAuth{}

\maketitle

\begin{abstract}
\section{}
\textbf{Introduction:} Public companies increasingly deploy and develop artificial-intelligence (AI) systems, and a growing public record documents when those systems fail; whether such failures reach investors through securities disclosure has not been measured. \textbf{Methods:} We link every incident in the AI Incident Database (AIID) dated 2019--2026 (n = 1{,}389) to the organizations alleged to have deployed or developed the implicated system, retain those resolving to a U.S.-listed issuer filing domestic periodic reports (N = 307 incidents across 21 issuers), and search each issuer's Forms 8-K, 10-K, and 10-Q over the twelve months after the incident. Each incident is classified into a four-tier taxonomy: disclosed-specific (T1), legal-proceedings-only (T2), generic-risk-language (T3), and no disclosure located (T4). \textbf{Results:} Incident-specific disclosure is rare: 4 of 307 incidents, 1.3\% (Clopper--Pearson 95\% CI 0.4--3.3); substantive disclosure that names the event or its legal consequence reaches 7 of 307, 2.3\% (0.9--4.6). The complement --- 97.7\% (95.4--99.1; issuer-clustered bootstrap 97.5\%, 93.9--99.6) --- forms a disclosure ``shadow'' dominated not by silence but by substitution: 90.6\% of incidents (T3) are absorbed into generic AI risk-factor language that never names the event. We detect no evidence that disclosure increases with harm severity (Monte-Carlo exact p = 0.61); in a pre-specified secondary analysis it differs by the issuer's role as deployer versus developer (exact p = 0.02). The pattern is robust to issuer concentration (the shadow is 92.5\% after dropping the five largest issuers), to the search window (6/12/18/24 months), and to right-censoring, and an independent full-filing audit of 40 incidents found no missed disclosures. \textbf{Discussion:} The seven substantive disclosures are concentrated among incidents with a financial-statement or formal legal consequence (booked charges of approximately \$304M and \$478M in the two clearest cases), a pattern consistent with an accounting-and-legal disclosure channel under a materiality-gated voluntary regime. We report disclosure incidence, not legal compliance, and make no claim that any company breached a disclosure obligation. As mandatory AI-incident reporting arrives, these results provide a pre-mandate baseline.

\tiny
\keyFont{ \section{Keywords:} artificial intelligence incidents, securities disclosure, SEC filings, AI governance, materiality, corporate transparency, risk-factor disclosure, reproducibility }
\end{abstract}

\section{Introduction}

Two literatures meet in this paper without previously having been connected. The first catalogs AI failures. Since 2020 the AI Incident Database (AIID) has collected and structured public reports of situations in which AI systems caused, or nearly caused, real-world harm, giving each incident an identifier, a date, and structured tags naming the organizations alleged to have deployed or developed the system \citep{McGregor2021}. Parallel efforts --- standardized incident definitions \citep{OECD2024}, proposals for a mandatory reporting regime \citep{DixonFrase2025}, and taxonomies of AI harms \citep{Weidinger2022} --- share a recurring premise: that AI harms are under-reported to any authority, and that their true incidence is unknown. That premise is asserted widely but has not been quantified against any concrete disclosure channel.

The second literature studies what companies tell investors about AI. The AI Disclosures Project shows that AI now saturates U.S. securities filings and that the tone is overwhelmingly promotional \citep{Strauss2025}; complementary work indexes prospective AI risk-factor language in Form 10-K filings \citep{UbertiBonaMarin2025} and documents AI ``washing'' --- strategic over-claiming --- as its mirror image \citep{Song2026}. All of this measures the disclosure that \emph{is} filed. None measures the disclosure that is \emph{not} filed: the gap between what is publicly known to have happened and what reaches investors.

This paper measures that gap directly. We ask a narrow, answerable question: when an AIID-recorded AI incident involves a U.S.-listed company, is the incident disclosed to investors in that company's securities filings? We call the complement of the answer --- the share of recorded incidents that never surface in filings --- the \emph{AI incident disclosure shadow}. We measure disclosure \emph{incidence}, not legal compliance: whether a disclosure was required is a materiality determination we do not adjudicate.

\textbf{Theoretical framing and hypotheses.} We interpret the setting through voluntary-disclosure theory. Classical models predict ``unraveling'' toward full disclosure \citep{Grossman1981,Milgrom1981}, but unraveling breaks down under proprietary and reputational disclosure costs and uncertainty about what managers know \citep{Dye1985,Verrecchia2001,HealyPalepu2001}, and U.S. securities disclosure is in any case gated by \emph{materiality} --- a substantial likelihood that a reasonable investor would find a fact important \citep{TSC1976}, with contingent events assessed by probability multiplied by magnitude \citep{Basic1988}. This frame yields three testable expectations. \emph{H1 (incidence)}: under a materiality-gated voluntary regime, incident-specific disclosure of AI failures is rare. \emph{H2 (channel)}: substantive disclosure concentrates where an incident becomes a financial-statement or formal legal event --- an accounting and legal channel --- rather than where it is merely severe or publicly salient. \emph{H3 (severity-independence)}: because harm to third parties is distinct from financial impact to the issuer, disclosure incidence does not increase with the severity of the harm.

Our contribution is fourfold. First, we provide the first incident-level measurement of AI-failure disclosure by U.S.-listed issuers, with a four-tier taxonomy that distinguishes naming an incident, disclosing its legal consequence, and absorbing it into boilerplate. Second, we show the shadow's shape is substitution, not silence. Third, we present evidence consistent with the accounting-and-legal channel of H2 and reconcile the severity-independence of H3. Fourth, we release the annotated incident-to-filing dataset, coding logs, a hashed pre-registration, and a reproducible computational capsule. Two policy developments make a pre-mandate baseline timely: a 2026 U.S. Securities and Exchange Commission (SEC) rulemaking petition seeking mandatory AI governance disclosure \citep{SECpetition2026}, and Article 73 of the EU AI Act, which obliges serious-incident reporting for high-risk systems \citep{EUAIAct2024}.

\section{Related work}

\textbf{AI in the firm and in capital markets.} A fast-growing literature studies how artificial intelligence is adopted inside firms and how it reaches capital markets. Recent work in this journal alone spans AI for stock-price forecasting \citep{Rohan2026}, reinforcement-learning portfolio management, AI-based financial services and robo-advisory adoption \citep{Verma2026}, AI-driven financial risk management \citep{Jiang2026}, and the impact of AI on audit quality and cost in banking \citep{Johri2026}; a parallel strand connects firms' AI adoption to capital-market outcomes, including corporate ESG performance measured from annual reports \citep{Shen2025} and stock-price crash risk \citep{YangCrash2026}, while a governance strand formalizes the ethical and risk-management frameworks under which responsible adoption is supposed to occur \citep{Madanchian2025}. This body of work overwhelmingly studies the \emph{upside} of AI adoption or the frameworks meant to govern it. What it does not do --- and what we do here --- is measure whether the \emph{downside}, the concrete failures that AI systems produce, is transmitted to investors through the disclosure channel.

\textbf{Materiality and disclosure regulation.} The materiality standard defines what could require disclosure \citep{TSC1976,Basic1988}; empirical work shows mandates bind only with enforcement and that disclosure regulation has real effects \citep{ChristensenHailLeuz2016,LeuzWysocki2016}. We contribute incident-level evidence of how rarely known AI events are surfaced, without adjudicating materiality for any event.

\textbf{Risk-factor text and voluntary disclosure.} Risk-factor disclosures are informative in aggregate \citep{Campbell2014}, move investor perceptions \citep{KravetMuslu2013}, and are more valuable when specific \citep{Hope2016}, with a large filings-text literature building on dictionary methods \citep{LoughranMcDonald2011} and on machine-extracted risk types \citep{BaoDatta2014}. Voluntary-disclosure theory explains when unraveling toward full disclosure breaks down \citep{Dye1985,Verrecchia2001,HealyPalepu2001,Grossman1981,Milgrom1981}, and its empirical arm documents that managers disclose bad news selectively and often with delay: firms with bad news disclose less and later \citep{KothariShuWysocki2009}, disclosure of impending bad news is strategic \citep{Skinner1994}, and the reporting environment as a whole is shaped by these asymmetric incentives \citep{Beyer2010}. We add the observation that the realized events behind AI risk factors are almost never named, and we separate outright withholding from the substitution of a specific event by generic language.

\textbf{Cybersecurity disclosure: the closest analog.} Cybersecurity is the nearest disclosure-regime precedent and supplies our external-truth-versus-filing design. Announced breaches carry measurable market consequences \citep{Campbell2003,KannanReesSridhar2007}, firms underreport cyber-attacks relative to what capital markets later reveal \citep{Amir2018}, the security content of filings responds to regulation \citep{Gordon2006,Gordon2010}, and disclosed information-security risk factors are only weakly aligned with the risks that subsequently materialize \citep{WangKannanUlmer2013}. That last finding --- a wedge between disclosed risk language and realized events --- is precisely the wedge we measure for AI, but with the matched-incident denominator the cyber literature cannot observe because it lacks an external census of breaches.

\textbf{AI harms, auditing, and near-miss reporting.} Prospective AI-harm taxonomies and audit frameworks describe what could go wrong and how to check \citep{Weidinger2022,Raji2020,Bommasani2021}, and a nascent ecosystem studies who audits AI systems and how \citep{CostanzaChock2022}, while the AIID and OECD work catalog what did go wrong \citep{McGregor2021,OECD2024}. Outside AI, safety research shows that incident reporting depends on confidential, non-punitive channels and that measured-versus-reported gaps can be an order of magnitude \citep{Barach2000,Classen2011}, and that reporting systems systematically under-capture events without the right incentives \citep{Macrae2016}. AI has no analogous investor-facing reporting channel; we quantify the resulting gap.

\section{Materials and methods}

\subsection{Data and the matchable population}
We use the complete AIID snapshot of 27 July 2026 (SHA-256 recorded in the archive manifest), a database dump with CSV exports \citep{McGregor2021}; it contains 1{,}597 numbered incidents, of which 1{,}389 carry a date in 2019--2026. Each record includes canonical entity slugs for the alleged deployer and developer and an incident date; a CSET AI Harm Taxonomy classification is present for 213 incidents in the snapshot. Among the 307 matched in-window incidents, 40 carry CSET severity information; usable CSET information takes precedence, with the study rubric (Section~\ref{sec:sev}) applied where CSET information is unavailable or unusable; severity strata are therefore indicative rather than authoritative. The denominator is \emph{AIID-recorded, media-visible} incidents, not a census of all AI incidents (Section~\ref{sec:limits}).

We resolve each in-window incident's deployer and developer slugs to issuers through a curated crosswalk with status LISTED (a U.S. domestic 8-K/10-K/10-Q filer), PARENT (a subsidiary or product mapped to its listed parent), FOREIGN (a U.S.-listed foreign private issuer), or DELISTED, and confidence HIGH or MED. Generic-category slugs (e.g., ``deepfake-technology-developers''), individuals, non-corporate actors, and privately held firms are not matched. Of the 1{,}389 window incidents, \textbf{307 (22.1\%) resolve to a U.S. domestic-listed issuer} (the primary population; Figure~\ref{fig:flow}); the non-matches comprise privately held AI developers (243; e.g., OpenAI, Anthropic), other named organizations not resolved to a listed U.S. issuer (400), generic or individual actors (398), foreign issuers (26), and delisted firms (15). Because the 400 unresolved named organizations may include additional listed firms absent from the curated crosswalk, 307 is a conservative lower bound on the matchable population. The 307 incidents span 21 issuers and are concentrated (Herfindahl index 0.18; Alphabet 86, Meta 81, Tesla 36, Microsoft 35). Denominator construction is validated in Section~\ref{sec:val}: a 137-row human-validation sheet --- all low-confidence and PARENT mappings plus a random sample of high-confidence LISTED mappings --- accompanies the dataset, and we separately audit a sample of the unresolved named organizations for missed listed issuers. Two issuer-identifier errors surfaced during that review and were corrected, evidence that the layer functions.

\subsection{Disclosure taxonomy and search protocol}
Each matched incident receives exactly one tier: \textbf{T1 (disclosed-specific)}, a window filing references the incident as such; \textbf{T2 (legal-proceedings-only)}, the legal or regulatory consequence appears without being framed as an AI incident; \textbf{T3 (generic-risk-language)}, in-window filings carry AI or algorithm risk-factor language of the same technology or harm family but do not name the event; and \textbf{T4 (no disclosure located)}. For each incident we search the window [incident date, +12 months], forms 8-K/10-K/10-Q, restricted to the issuer's Central Index Key, using three term families --- incident-specific terms, legal-proceedings vocabulary, and generic AI terms --- via SEC EDGAR full-text search, and we hand-read the Risk Factors and Legal Proceedings sections of the first annual report following each incident. The three term families are applied in sequence: incident-specific terms (the product, victim, place, or docket associated with the event), legal-proceedings vocabulary (the counterparty name, regulator, or matter caption), and generic AI terms (a fixed lexicon of AI, machine-learning, algorithm, and automation risk language). A verdict of T4 therefore requires that all three families return nothing in the window, not merely that the incident is unnamed. All queries are logged as deterministic URLs so every verdict is reproducible. Full coding rules and a decision tree are provided in the released codebook and pre-registration, and the retrieval protocol is validated independently in Section~\ref{sec:val}.

\subsection{Use of AI assistance}
Entity-resolution candidates, per-incident disclosure searches, and analysis code were produced with assistance from a large-language-model system (Anthropic Claude, Opus-class model, accessed August 2026 via the Claude Agent SDK) under human validation, which we disclose in the interest of research integrity (see also Acknowledgments). The system \emph{proposed} labels from deterministic, logged EDGAR queries and filing quotations; the codebook and decision rules were human-specified; every label is traceable to its search record; and an independent second coding pass was performed blind to the first. Crucially, the possibility that AI-assisted retrieval missed a disclosure is not left as an unmeasured dependency: it is assessed independently in Section~\ref{sec:val} by a search-free human reading of complete filings, which found no missed disclosures. All statistics are generated by deterministic scripts (fixed seed 42) and independently re-derived by a verifier. No AI system is listed as an author.

\subsection{Coding, reliability, and statistics}
Pass 1 coded all 307 incidents from the search records; Pass 2 independently re-coded a stratified reliability subsample (n = 68: 28 of the 29 non-T3 incidents --- all 4 T1, all 3 T2, and 21 of the 22 T4 --- plus a seed-42 random 40 of the 278 T3 incidents) directly from raw filings. Because T3 dominates, we report Cohen's $\kappa$ alongside the prevalence-robust PABAK and Gwet's AC1 \citep{Cohen1960,Gwet2008}, per-tier agreement, and an inverse-probability-weighted projection to the full sample. Proportions carry both Wilson \citep{Wilson1927} and, given the small number of positives, exact Clopper--Pearson 95\% intervals \citep{ClopperPearson1934}. Because incidents from the same issuer share disclosure language and are not independent, we also report an issuer-clustered bootstrap (10{,}000 resamples, seed 42); because it rests on only 21 issuer clusters, its precision is limited, which is why we present it alongside --- not in place of --- the exact binomial intervals and the concentration checks below. To test whether substantive disclosure differs across the severity strata (the H3 test) and, as a pre-specified secondary analysis, by issuer role, we use Monte-Carlo exact permutation tests (20{,}000 permutations) rather than asymptotic $\chi^2$, appropriate under sparse positives. A Firth-penalized logistic specification \citep{Firth1993} is identified as future work once a larger positive count is available; with seven substantive positives we deliberately do not present a multivariable model as confirmatory. Concentration robustness comprises leave-one-issuer-out and drop-top-$k$ recomputations; window robustness re-codes at 6/18/24 months and restricts to complete (non-right-censored) windows. The analysis plan was hashed before coding; every number is script-generated and re-derived by an independent verifier, and is byte-identical across runs.

\subsection{Severity coding}\label{sec:sev}
Because H3 turns on severity, we give the variable the same transparency as the disclosure taxonomy. CSET AI Harm Taxonomy severity information is available for 40 of the 307 matched incidents and takes precedence where usable. Incidents without usable CSET severity information are coded to a three-level rubric aligned to CSET concepts: \textbf{Severe} --- death, serious physical injury, or large-scale or systemic violation of fundamental rights (e.g., an autonomous-vehicle collision, a safety-critical failure, or mass surveillance); \textbf{Moderate} --- non-trivial financial, legal, or discriminatory harm to identifiable parties short of serious physical injury; and \textbf{Limited} --- isolated, reversible, or primarily reputational harm, including most content and misinformation incidents without demonstrated downstream physical or financial harm. Incidents with multiple harms are coded to the most severe applicable level. The severity coding yields Severe = 42, Moderate = 50, and Limited = 215. Two properties bound its use: severity is single-pass (not independently double-coded like the disclosure tiers), and CSET and rubric-coded observations are combined within each stratum, so \emph{the severe tier (42) is not identical to the CSET-classified subset (40)} and severity strata are read as indicative rather than authoritative. The 40 CSET-classified incidents, their recorded CSET values, and the assigned study severity are listed in Table~\ref{tab:cset}; all rubric criteria and per-incident severity codes are released with the dataset.

\subsection{Retrieval and entity-resolution validation}\label{sec:val}
Two independent audits bound the study's principal validity threats.

\emph{Retrieval (false negatives).} To test whether the search-based protocol missed disclosures, a random 40 incidents --- 20 coded T3 and 20 coded T4, drawn seed-42 within tier --- were re-examined by reading \emph{every} in-window 8-K, 10-K, and 10-Q for the issuer in full, independently of the search results that produced the original verdict. This search-free re-read found zero missed substantive disclosures. With 0 misses in 40, the exact one-sided 95\% upper bound on the per-incident miss rate is 7.2\% (equivalently, the rule of three gives $\approx$7.5\%); reading 60 or 80 incidents would tighten this to $\approx$4.9\% or $\approx$3.7\% and is the natural way to strengthen the bound further. The direction of any residual error is important: a missed disclosure would reclassify a genuine T1/T2 incident as T3/T4, which would bias the estimated shadow \emph{upward}, toward the headline --- which is precisely why an independent, search-free audit is the right check. Separately, all seven substantive (T1/T2) incidents were re-examined directly in the underlying SEC filings, not inferred from the sampling frame.

\emph{Entity resolution (denominator).} The matchable population depends on the crosswalk. Beyond the 137-row validation sheet (all low-confidence and PARENT mappings plus a random sample of high-confidence LISTED mappings, which surfaced and corrected two identifier errors), we reviewed the highest-frequency unresolved named organizations for missed U.S.-listed issuers. They are dominated by disinformation campaigns (e.g., Storm-1516, Spamouflage, Doppelganger, Operation Overload), named individuals, and government or law-enforcement agencies, with the only company-like entities being privately held or foreign; the review surfaced no additional U.S.-listed issuer. We therefore treat 307 as a conservative estimate of the matchable population while noting that crosswalk omissions among the long tail of low-frequency unresolved organizations cannot be fully excluded.

\section{Results}

\subsection{The disclosure shadow (H1)}
Across 307 incidents the distribution is T1 = 4, T2 = 3, T3 = 278, T4 = 22 (Table~\ref{tab:dist}). The directly observed rates are the primary quantities: incident-specific disclosure is 4 of 307, 1.3\% (Clopper--Pearson 0.4--3.3), and substantive disclosure (T1+T2) is 7 of 307, 2.3\% (0.9--4.6). Their complement is the shadow (T3+T4), 97.7\% (95.4--99.1; issuer-clustered bootstrap 97.5\%, 93.9--99.6). The shadow is dominated by \emph{substitution}: nine in ten incidents fall in T3, where the issuer carries same-family AI risk language but never references the incident. The rarity of incident-specific disclosure is consistent with H1.

\subsection{Severity, role, and concentration}
We detect no evidence that disclosure incidence rises with severity (Table~\ref{tab:sev}): the Monte-Carlo exact permutation test on substantive disclosure across severity tiers gives p = 0.61 --- no detectable severity gradient, consistent with H3 (a null test is not proof of exact independence). In a pre-specified secondary analysis (role is not among H1--H3), disclosure differs by the issuer's role (exact p = 0.02): all four T1 incidents involve an issuer acting as deployer, and no incident in which the issuer's only alleged role is developer (n = 58) is specifically disclosed --- an association consistent with disclosure tracking the operational and financial exposure that accrues to a deployer. The pattern is not an artifact of a few firms: the top two issuers supply 54\% of incidents, yet the shadow is 96.4\% after dropping them and 92.5\% after dropping the top five (Table~\ref{tab:rob}; Figure~\ref{fig:pareto}), and leave-one-issuer-out shadow rates range only 96.8--98.6\%. Only 4 of 21 issuers have any substantive disclosure.

\subsection{The accounting and legal channel (H2)}
The seven substantively disclosed incidents illustrate the pattern (Table~\ref{tab:acct}). Two carry a booked financial-statement impact, and these are the two clearest specific disclosures: a real-estate group's wind-down of an algorithmic home-buying business, disclosed with an inventory write-down of approximately \$304 million, and an automotive issuer's disclosure of an autonomous-vehicle subsidiary's pause after a pedestrian incident, narrated alongside a \$478 million restructuring charge. The remainder are litigation or regulatory matters surfacing in Legal Proceedings or Contingencies notes. By incident category, substantive disclosure is concentrated among autonomous-vehicle, copyright, and employment incidents, while the two largest categories --- misinformation/content (n = 70) and generative-AI errors (n = 57) --- contain zero substantive disclosures. This is consistent with H2 (an association, not a formally identified mechanism): substantive disclosure is concentrated among incidents with a financial-statement or formal legal consequence, and we detect no association with severity. The boundary between substitution (T3) and silence (T4) tracks the year each issuer first added AI risk-factor language to its 10-K (Table~\ref{tab:onset}).

\subsection{Temporal, category, and issuer structure}
The shadow does not fade as AI disclosure practice matures. Across the eight incident years the shadow rate never falls below 93.6\%, and we detect no linear downward trend in substantive disclosure (Cochran--Armitage z = $-0.28$, p = 0.78; Figure~\ref{fig:year}). The four specific disclosures appear only in 2021 and 2023 and coincide with booked charges rather than with any general improvement in AI transparency, even as the share of issuers carrying AI risk-factor language rose steadily over the period (Table~\ref{tab:onset}). By harm category the concentration of substantive disclosure is stark: it appears only in autonomous-vehicle, copyright/IP, employment, and a single discrimination incident, whereas the two largest categories --- misinformation and content (n = 70) and generative-AI errors (n = 57) --- together contribute 127 incidents and \emph{zero} substantive disclosures. The pattern is also not driven by the most-scrutinized firms: the two highest-volume issuers, which account for 54\% of all incidents, disclose essentially nothing specific (shadow rates of 100\% and 98\%), while the only issuer with more than one specific disclosure is an automotive deployer whose incidents produced booked charges and a federal safety proceeding. Substantive disclosure therefore tracks the financial and legal consequence of an incident, not the volume of a firm's incidents or the public salience of the harm.

\subsection{Reliability and robustness}
On the double-coded subsample, agreement is 85.3\% and Cohen's $\kappa = 0.73$; the prevalence-robust PABAK = 0.71 and Gwet's AC1 = 0.82 are higher, and the inverse-probability-weighted projection is 90.9\%. The full Pass-1 $\times$ Pass-2 confusion matrix and per-tier agreement are released with the archive; disagreements concentrate at the T2/T3 boundary, and four of the ten disagreements cross the substantive/non-substantive line, bounding the sensitivity of the 2.3\% substantive rate. The headline is stable under every pre-registered perturbation (Table~\ref{tab:rob}): excluding pre-2021 incidents gives 97.3\%; the search window gives 99.0/97.7/97.7/97.7\% at 6/12/18/24 months; restricting to the 252 complete (non-right-censored) windows gives 97.2\% (94.4--98.9); and the full-filing re-audit of 40 incidents found zero missed disclosures (exact one-sided 95\% upper bound on the miss rate 7.2\%; Section~\ref{sec:val}). Crucially, the shadow persists among plausibly material incidents: among the 42 severe-tier incidents it is 95.2\% (two substantive), and among the 254 incidents at the five largest, most-scrutinized issuers it is 98.8\%.

\section{Discussion}

\subsection{What the shadow is, and is not}
AIID-recorded AI incidents involving U.S.-listed companies are almost never disclosed as incidents to investors. But the shadow is mostly not silence; it is the substitution of a specific event by generic risk language --- firms disclose the \emph{category} of harm their AI could cause while omitting the \emph{instances} that occurred. This is a subtler transparency gap than non-disclosure, and complements, from the opposite direction, the finding that filed AI disclosures are overwhelmingly promotional \citep{Strauss2025}. We emphasize that ``no related disclosure located'' is a search-scoped measure of disclosure \emph{incidence}, not a finding of non-compliance: U.S. disclosure is materiality-gated, and many individual incidents may not have met the materiality threshold. The contribution is to measure the incidence gap that the literature has asserted but never quantified.

T3 is not a mechanical residual assigned to every issuer that had an incident. An incident is coded T3 only when the issuer's in-window filings already carry same-family AI risk-factor language; incidents that occur before an issuer first adopts such language fall into T4, which is why T4 is 7.2\% rather than near zero. This timing structure --- issuers adopt AI risk-factor language at different fiscal years --- is visible directly in Table~\ref{tab:onset}, and it is what makes the T3/T4 split informative rather than definitional.

\subsection{Interpretation: the accounting and legal channel}
The concentration of substantive disclosure among autonomous-vehicle, copyright, and employment incidents, the two booked charges among seven substantive cases, and the absence of a severity gradient are together consistent with an accounting-and-legal channel (H2). Voluntary-disclosure theory predicts this pattern: absent a mandate, and with proprietary and reputational costs of naming a failure, a rational issuer discloses when an incident becomes a booked charge or a reportable proceeding that a reasonable investor would find material \citep{Dye1985,Verrecchia2001,Basic1988}. The absence of a severity gradient (H3) is what this account anticipates: harm to third parties (the severity axis) is distinct from financial impact to the issuer (the materiality axis), so incidence should track the latter, not the former. We frame the channel as an interpretation consistent with the evidence rather than a formally identified mechanism: with seven substantive positives, a logistic model of disclosure on a booked-impact indicator with issuer-clustered errors is underpowered, and we flag it as the natural next step on a larger sample.

\subsection{Why substitution, not silence}
The dominant tier is not silence (T4) but substitution (T3): firms disclose the category of AI harm while omitting the instances that occurred. Three forces plausibly combine to produce this. First, materiality gating means most single incidents --- absent a booked charge or a formal proceeding --- fall below the threshold at which a reasonable investor's decisions would change, so nothing specific is required \citep{TSC1976,Basic1988}. Second, when disclosure is discretionary, the proprietary and reputational cost of naming one's own failure is real, and voluntary-disclosure theory predicts selective, delayed release of bad news exactly under these conditions \citep{Dye1985,Verrecchia2001,Skinner1994,KothariShuWysocki2009}. Third, generic risk-factor language is close to costless and may reduce the informational significance of omitting a specific event: a firm that has already stated that its AI ``may produce inaccurate outputs'' provides broader ex ante risk disclosure, weakening any inference an outside reader would draw from the omission of a specific event. The result is an equilibrium in which the boilerplate expands to blanket the technology while the incidents behind it stay unnamed --- the same wedge between disclosed risk language and realized events documented for information-security risk factors \citep{WangKannanUlmer2013}, and a caution that text-based measures of AI risk built from filings \citep{BaoDatta2014} capture the stated category, not the realized incident. Substitution, on this reading, is not an accident of drafting but the predictable output of a materiality-gated, cost-sensitive disclosure regime.

\subsection{Policy and practical implications}
Two developments would create the incident-routing mechanism our results show is currently absent: the 2026 SEC rulemaking petition \citep{SECpetition2026} and Article 73 of the EU AI Act \citep{EUAIAct2024}. This paper supplies the pre-intervention baseline against which either regime's effect could be measured. For investors and auditors, the T3-dominant pattern is a caution that AI risk-factor language is a weak signal of incident exposure, and a firm's ratio of specific to generic disclosure could serve as a diligence flag. For standard-setters and incident databases, adding issuer identifiers to catalogs such as AIID would make this measurement routine. The cybersecurity Item 1.05 regime (29 material and 50 voluntary filings in its first two years) is best read as institutional context, not a comparable rate: it counts disclosures filed under a mandate and cannot observe the cyber shadow, whereas our design supplies the matched-incident denominator \citep{Amir2018,SEC2023cyber}.

\subsection{Limitations}\label{sec:limits}
Several bound the interpretation. (i) The denominator is a curated, media-derived catalog, not a census; the matchable population is issuer-concentrated, and 307 is a lower bound because some unresolved named organizations may be listed. (ii) Parent/product mapping required judgment; two identifier errors were caught and corrected, and residual errors cannot be excluded. (iii) Detection relied on full-text search plus targeted reads; the 40-incident audit found no misses, but very obliquely worded disclosures could be missed, biasing the specific rate downward. (iv) The 12-month window is a choice; results are robust to 18/24 months and to restricting to complete windows, and mildly higher at 6 months. (v) CSET severity covers 40 of 307 incidents; the rest are rubric-coded, so severity strata are indicative. (vi) Primary codes are single-pass with an independent reliability subsample; the T2/T3 boundary carries the most residual uncertainty. (vii) Inference rests on few positive events; we therefore report exact intervals and permutation tests, but statistical power for stratum comparisons is limited.

These bounds also mark the natural extensions. A larger matchable population --- obtained by adding issuer identifiers to incident catalogs, or by pooling AIID with complementary databases --- would raise the positive count enough to fit the issuer-clustered logistic model of disclosure on a booked-impact indicator that H2 motivates but our seven positives cannot support. Linking incidents to filing dates would allow an event-study of whether the rare specific disclosures move prices, testing materiality directly rather than by proxy. Cross-jurisdictional replication under the EU AI Act's Article 73 regime would convert our pre-mandate baseline into a difference-in-differences design once serious-incident reporting begins. Finally, the same matched-incident method could be turned on private-firm channels (venture disclosures, litigation dockets) to characterize the shadow where no periodic-reporting obligation exists at all.

\subsection{Generalizability}
The design generalizes beyond this snapshot in two directions. Methodologically, external-truth-versus-filing matching is portable to any domain with an independent event census and a disclosure channel; cybersecurity is the established analog \citep{Amir2018,WangKannanUlmer2013}, and product-safety and environmental-incident registries are natural next targets. Substantively, the direction of selection plausibly makes the estimated shadow conservative rather than inflated: our denominator is limited to media-visible, AIID-recorded incidents at large, heavily scrutinized issuers, precisely the incidents most likely to be disclosed. We therefore expect that broadening the census toward less-visible incidents would, if anything, widen the gap; but because the disclosure behavior of unrecorded incidents is unobservable, this direction cannot be established empirically and we do not claim the headline itself is a formal lower bound.

\subsection{Conclusion}
When an AIID-recorded AI incident involves a U.S.-listed company, the incident is almost never disclosed to investors as such; the gap is large and structured, and the rare specific disclosures are concentrated among incidents with financial-statement or legal consequences, while we detect no association with severity. As mandatory AI-incident reporting arrives on both sides of the Atlantic, this measurement provides its missing baseline. The study characterizes disclosure behavior under the current voluntary regime and makes no determination about any company's legal obligations.

\section*{Conflict of Interest Statement}
The authors declare that the research was conducted in the absence of any commercial or financial relationships that could be construed as a potential conflict of interest.

\section*{Author Contributions}
SC conceived and designed the study, curated the data, performed the analysis, and wrote the manuscript. RC contributed to the study design, the coding protocol and independent verification, and the critical revision of the manuscript. Both authors approved the final version.

\section*{Funding}
This research received no external funding.

\section*{Acknowledgments}
Analysis and drafting were performed with assistance from a large-language-model system (Anthropic Claude, Opus-class model, accessed August 2026 via the Claude Agent SDK) used for entity-resolution candidate generation, disclosure-search execution, and analysis code, in all cases under human validation and with full logging; the corresponding author is responsible for the accuracy of all content, citations, and references (see Materials and Methods, ``Use of AI assistance'').

\section*{Ethics Statement}
This study did not involve human participants or animals. It uses only publicly available records (the AI Incident Database and SEC EDGAR). The study measures disclosure behavior in the aggregate and makes no assertion that any named issuer violated a disclosure duty.

\section*{Data Availability Statement}
The annotated incident-to-issuer match table, per-incident search records and coding logs, the independent-coding disagreement log, the hashed pre-registration, and all analysis and verification scripts are openly available at \url{https://github.com/samirrc2/ai-incident-disclosure-shadow} and archived on Zenodo (DOI to be inserted on acceptance). A fully offline, deterministic computational capsule reproduces every reported number and figure. The underlying AIID snapshot is publicly available from the Responsible AI Collaborative; the pinned version and its SHA-256 are recorded in the manifest.

\bibliographystyle{Frontiers-Harvard}
\bibliography{references}

%%%%%%%%%%%%%%%%%%%%%%%%%%%%%%%%%%%%%%%%%%%%%%%%%  TABLES (Frontiers: at end)  %%%%%%%%%%%%%%%%%%%%%%%%%%%%%%%%%%%%%%%%%%%%%%%%%

\begin{table}[h!]
\caption{Four-tier disclosure distribution (N = 307). Intervals are Clopper--Pearson exact 95\%.}
\label{tab:dist}
\centering
\begin{tabular}{@{}llrl@{}}
\toprule
Tier & Meaning & n & Share (95\% CI) \\ \midrule
T1 & Disclosed-specific        & 4   & 1.3\% (0.4--3.3) \\
T2 & Legal-proceedings-only    & 3   & 1.0\% \\
T3 & Generic-risk-language     & 278 & 90.6\% (86.7--93.4) \\
T4 & No disclosure located     & 22  & 7.2\% (4.6--10.7) \\ \midrule
\multicolumn{2}{@{}l}{Shadow (T3+T4)} & 300 & 97.7\% (95.4--99.1) \\
\bottomrule
\end{tabular}
\end{table}

\begin{table}[h!]
\caption{Disclosure by incident severity tier. Monte-Carlo exact permutation test on substantive disclosure across tiers: p = 0.61.}
\label{tab:sev}
\centering
\begin{tabular}{@{}lrrrr@{}}
\toprule
Severity & n & T1 specific & T1+T2 & Shadow (T3+T4) \\ \midrule
Severe   & 42  & 2.4\% & 4.8\% & 95.2\% \\
Moderate & 50  & 0.0\% & 2.0\% & 98.0\% \\
Limited  & 215 & 1.4\% & 1.9\% & 98.1\% \\ \bottomrule
\end{tabular}
\end{table}

\begin{table}[h!]
\caption{All substantively disclosed incidents (T1/T2) and their disclosure channel.}
\label{tab:acct}
\centering
\begin{tabular}{@{}llll@{}}
\toprule
Issuer & Tier & Filing location & Booked impact \\ \midrule
Zillow Group    & T1 & 8-K Item 8.01                 & Yes --- $\approx$\$304M write-down \\
General Motors  & T1 & 10-K Item 1/1A                & Yes --- \$478M restructuring \\
General Motors  & T1 & 10-K Item 1/1A                & No incident-specific reserve \\
Meta Platforms  & T1 & Legal Proceedings (named suit) & No \\
Meta Platforms  & T2 & 10-K Contingencies           & No \\
Tesla           & T2 & Legal Proceedings / regulatory & No \\
General Motors  & T2 & 10-Q Contingencies           & No \\ \bottomrule
\end{tabular}
\end{table}

\begin{table}[h!]
\caption{Robustness of the shadow rate (T3+T4). CIs are issuer-clustered bootstrap or Clopper--Pearson where shown.}
\label{tab:rob}
\centering
\begin{tabular}{@{}lrl@{}}
\toprule
Specification & n & Shadow \\ \midrule
Full primary sample                     & 307 & 97.7\% (95.4--99.1) \\
Issuer-clustered bootstrap              & 307 & 97.5\% (93.9--99.6) \\
Exclude pre-2021 incidents              & 257 & 97.3\% \\
Drop top-2 issuers                      & 140 & 96.4\% \\
Drop top-5 issuers                      & 53  & 92.5\% \\
Complete (non-right-censored) windows   & 252 & 97.2\% (94.4--98.9) \\
6-month window                          & 307 & 99.0\% \\
12-month window (primary)               & 307 & 97.7\% \\
18-month window                         & 307 & 97.7\% \\
24-month window                         & 307 & 97.7\% \\
Severe-tier subset                      & 42  & 95.2\% \\
Top-5-issuer subset                     & 254 & 98.8\% \\ \bottomrule
\end{tabular}
\end{table}

\begin{table}[h!]
\caption{Fiscal year in which each issuer first added AI risk-factor language to its Form 10-K --- the T3/T4 boundary. An incident dated before its issuer's onset year falls into T4 (no disclosure located); one dated on or after falls into T3 (generic risk language). Issuers without a determinable onset year are omitted.}
\label{tab:onset}
\centering
\begin{tabular}{@{}lc@{}}
\toprule
Issuer & First fiscal year with AI risk-factor language in 10-K \\ \midrule
Alphabet       & 2018 \\
Amazon         & 2019 \\
Meta           & 2019 \\
Tesla          & 2019 \\
Apple          & 2023 \\
Microsoft      & 2023 \\
General Motors & 2024 \\ \bottomrule
\end{tabular}
\end{table}

{\footnotesize
\setlength{\LTcapwidth}{\textwidth}
\begin{longtable}{@{}r p{9.4cm} l@{}}
\caption{CSET-classified incidents and study severity: the 40 of 307 matched incidents for which CSET AI Harm Taxonomy severity is available. The two CSET severity fields are reproduced verbatim from the source data.}\label{tab:cset}\\
\toprule
AIID ID & CSET classification (harm level; tangible harm) & Study severity \\ \midrule
\endfirsthead
\multicolumn{3}{@{}l}{\emph{Table \ref{tab:cset} continued}}\\ \toprule
AIID ID & CSET classification (harm level; tangible harm) & Study severity \\ \midrule
\endhead
\midrule \multicolumn{3}{r@{}}{\emph{continued on next page}}\\
\endfoot
\bottomrule
\endlastfoot
81 & none; no tangible harm, near-miss, or issue & Moderate \\
82 & none; no tangible harm, near-miss, or issue & Moderate \\
83 & none; no tangible harm, near-miss, or issue & Moderate \\
84 & none; no tangible harm, near-miss, or issue & Moderate \\
89 & none; tangible harm definitively occurred & Severe \\
92 & AI tangible harm event; tangible harm definitively occurred & Severe \\
97 & AI tangible harm issue; non-imminent risk of tangible harm (an issue) occurred & Moderate \\
102 & none; no tangible harm, near-miss, or issue & Moderate \\
105 & AI tangible harm event; tangible harm definitively occurred & Severe \\
113 & none; no tangible harm, near-miss, or issue & Moderate \\
116 & AI tangible harm event; tangible harm definitively occurred & Severe \\
124 & AI tangible harm event; tangible harm definitively occurred & Severe \\
125 & AI tangible harm event; tangible harm definitively occurred & Severe \\
127 & none; tangible harm definitively occurred & Severe \\
129 & none; no tangible harm, near-miss, or issue & Moderate \\
139 & none; no tangible harm, near-miss, or issue & Moderate \\
142 & AI tangible harm issue; non-imminent risk of tangible harm (an issue) occurred & Moderate \\
143 & none; no tangible harm, near-miss, or issue & Moderate \\
144 & unclear; unclear & Limited \\
145 & AI tangible harm issue; non-imminent risk of tangible harm (an issue) occurred & Moderate \\
149 & AI tangible harm event; tangible harm definitively occurred & Severe \\
153 & AI tangible harm event; tangible harm definitively occurred & Severe \\
156 & none; tangible harm definitively occurred & Severe \\
157 & unclear; tangible harm definitively occurred & Severe \\
160 & AI tangible harm issue; non-imminent risk of tangible harm (an issue) occurred & Moderate \\
220 & AI tangible harm event; tangible harm definitively occurred & Severe \\
281 & none; non-imminent risk of tangible harm (an issue) occurred & Moderate \\
313 & none; no tangible harm, near-miss, or issue & Moderate \\
354 & AI tangible harm event; tangible harm definitively occurred & Severe \\
360 & none; no tangible harm, near-miss, or issue & Moderate \\
367 & none; no tangible harm, near-miss, or issue & Moderate \\
393 & none; no tangible harm, near-miss, or issue & Moderate \\
414 & none; no tangible harm, near-miss, or issue & Moderate \\
489 & none; no tangible harm, near-miss, or issue & Moderate \\
501 & AI tangible harm event; tangible harm definitively occurred & Severe \\
554 & none; no tangible harm, near-miss, or issue & Moderate \\
574 & none; no tangible harm, near-miss, or issue & Moderate \\
576 & none; no tangible harm, near-miss, or issue & Moderate \\
583 & none; no tangible harm, near-miss, or issue & Moderate \\
614 & none; no tangible harm, near-miss, or issue & Moderate \\
\end{longtable}
}

\noindent \emph{Note:} CSET severity information was available for 40 of the 307 matched incidents. The two CSET fields are reproduced verbatim above, including source-recorded ``unclear'' values. CSET information took precedence where usable; otherwise severity was assigned using the rubric described in Section~\ref{sec:sev}. The remaining 267 incidents were rubric-coded. Complete incident-level codes are available in the released dataset.


%%%%%%%%%%%%%%%%%%%%%%%%%%%%%%%%%%%%%%%%%%%%%%%%%  FIGURE CAPTIONS + FIGURES (Frontiers: at end)  %%%%%%%%%%%%%%%%%%%%%%%%%%%%

\begin{figure}[h!]
\begin{center}\includegraphics[width=0.82\textwidth]{figures/fig1.png}\end{center}
\caption{Sample construction: from 1{,}597 AIID incidents to the primary matchable population of N = 307 (21 issuers), with exclusion branches.}\label{fig:flow}
\end{figure}

\begin{figure}[h!]
\begin{center}\includegraphics[width=0.9\textwidth]{figures/fig2.png}\end{center}
\caption{Disclosure tier by incident year (share of incidents).}\label{fig:year}
\end{figure}

\begin{figure}[h!]
\begin{center}\includegraphics[width=0.62\textwidth]{figures/fig3.png}\end{center}
\caption{Shadow rate (T3+T4) by incident severity, with 95\% Wilson intervals.}\label{fig:sev}
\end{figure}

\begin{figure}[h!]
\begin{center}\includegraphics[width=0.9\textwidth]{figures/fig4.png}\end{center}
\caption{Incident concentration by issuer (Pareto): incidents per issuer and cumulative share.}\label{fig:pareto}
\end{figure}


\end{document}
